\subsection{可用储能容量的全年消融}\label{sec:capacity}
可用库存空间既能把低价电转移到高价时段，也能为预报误差留出充放电余地。为量化这部分物理资源的作用，保持真实初末库存6,000 kWh、5,000 kW功率和每向0.9效率，取可用窗口宽度$W\in\{0,4800,9600\}$ kWh：
\begin{equation}\label{eq:capacitywindow}
 \mathcal E_W=[6000-W/2,\,6000+W/2],\qquad E_0=E_T=6000.
\end{equation}
三档窗口分别为$[6000,6000]$、$[3600,8400]$和$[1200,10800]$ kWh。$W=0$表示停用库存调节，单向执行下充放电均为零；$W=9600$对应主模型的可用容量。该设计比较同一储能设备开放的运行空间，保持初始能量与功率不变。

\pp{10}{A：扩大可用窗口的模型关系}固定外生过程分布、信息与交易规则，令$\Pi(W)$为窗口$\mathcal E_W$内全部合法因果策略，$J^*(W)$为其期望费用下确界。则
\begin{equation}\label{eq:capacitynesting}
 W_1\le W_2\ \Longrightarrow\ \Pi(W_1)\subseteq\Pi(W_2)
 \ \Longrightarrow\ J^*(W_2)\le J^*(W_1).
\end{equation}
\begin{proof}
区间以同一库存为中心，故$\mathcal E_{W_1}\subseteq\mathcal E_{W_2}$。窄窗口策略的相同合同与电池动作在宽窗口中仍满足状态、功率、信息和终点约束；目标不变，对可行策略集合取下确界即得结论。外网追索与弃电保证保持6,000 kWh库存的策略可行。
\end{proof}

实验对四种机制分别连续回放334日，每一窗口重新规划合同并执行电池控制。场景、视域、预报权限及末值规则保持主配置；半窗口使用161节点，全窗口321节点，间距同为30 kWh。零窗口采用单点状态分支。规划器、值函数和执行器使用相同窗口边界，实际账单按原结算式重算。
\figfixed{storage-capacity}{三档可用库存窗口的实际费用：相同初末库存和功率\label{fig:capacity}}{1.0}
\textbf{容量作用：}在问题二、问题三、问题四日前和日内机制中，完整9.60 MWh窗口相对停用储能，费用依次节省496.61万元、节省518.01万元、节省487.01万元、节省509.77万元。将可用窗口从4.80 MWh扩大到9.60 MWh，费用依次节省172.72万元、节省176.23万元、节省173.10万元、节省176.45万元。完整费用、分项账目及配对日差区间见附录表\ref{tab:capacitycost}至表\ref{tab:capacityci}。容量对照量化储能可用空间的费用作用；主表固定9.60 MWh窗口的跨日贪心对照，则量化控制策略的改进。

命题10针对精确模型，近似闭环的费用排序由完整回放确定。图中连线连接三个已计算容量点，不据此推断未计算容量的费用或一般边际收益递减关系。费用改善衡量可用库存空间在本控制框架中的运行价值，额定容量配置还需另行纳入功率配比和投资成本。
